\section{Related Work}
\label{sec:related}

\subsection{Constraint Handling in Process Control}

MPC is the dominant constrained control methodology in the process industries~\cite{qin2003survey,mayne2000constrained}. By solving a constrained optimization over a receding horizon, MPC explicitly handles input and state constraints. Offset-free MPC~\cite{morari1999model} augments the state with disturbance estimates to achieve zero steady-state offset under plant-model mismatch. Robust MPC formulations---including tube MPC, min-max MPC, and stochastic MPC with chance constraints---provide constraint satisfaction guarantees under bounded or stochastic uncertainty, but at the cost of increased conservatism and computational complexity. Cautious MPC~\cite{hewing2020cautious} incorporates GP uncertainty into the prediction model, propagating posterior variance through the horizon to construct chance constraints. However, the computational cost of GP-augmented MPC limits its real-time applicability for fast process control loops, and the safety guarantee is horizon-dependent rather than providing per-step forward invariance.

Economic MPC optimizes process economics directly rather than tracking a setpoint, and process safety constraints are typically enforced as hard output constraints within the MPC formulation. Disturbance-observer-based control provides an alternative approach to handling plant-model mismatch without full MPC optimization, using estimated disturbances for feedforward compensation. These approaches are complementary to the safety filter proposed here: the Robust HOCBF filter can be deployed alongside any upstream controller (MPC, LQR, PID, or a learned policy) to provide an independent, per-step safety guarantee that does not depend on prediction horizon length or disturbance model accuracy.

\subsection{Barrier-Function-Based Safety Filtering}

CBFs~\cite{ames2017cbf,wieland2007constructive} enforce forward invariance of a safe set through a QP that minimally modifies the control action. For constraints with relative degree $m \geq 2$, HOCBFs~\cite{xiao2021hocbf} construct a recursive $\psi$-chain of class-$\mathcal{K}$ functions. This is particularly relevant for process control: steam pressure, temperature, and composition constraints often have relative degree $m \geq 2$ with respect to available control inputs (valve positions, feedwater flow, fuel feed rate). Robust CBF formulations~\cite{taylor2020issf,jankovic2018robust} add constant or state-dependent margins based on worst-case disturbance bounds. Adaptive CBFs~\cite{long2023ecbf} estimate uncertain parameters online. Cosner et al.~\cite{cosner2023robust} formulate input-to-state safe CBFs with structural disturbance bounds. These approaches rely on a priori disturbance characterizations rather than data-driven, calibrated uncertainty estimates.

\subsection{GP-Based Uncertainty Modelling for Process Control}

GPs provide a nonparametric Bayesian framework for learning unknown functions from data with calibrated uncertainty~\cite{rasmussen2006gp}. In process control, GPs have been applied to soft sensing, model predictive control, and fault detection. The PAC-Bayes concentration inequality~\cite{srinivas2010gp} provides a high-probability bound on GP prediction error, enabling safety-critical applications. Several works combine GPs with CBFs~\cite{cheng2019guaranteed,fisac2018general,castaneda2021pointwise}, but these apply GP uncertainty only at the top-level CBF constraint ($m=1$) or over an MPC horizon, without compositional propagation through the HOCBF $\psi$-chain. For process constraints with $m \geq 2$, this omission is structurally significant: the model residual propagates through successive Lie derivatives and class-$\mathcal{K}$ function compositions, and a single-level bound cannot capture the amplification of uncertainty through the chain. Lederer et al.~\cite{lederer2021uniform} derive uniform error bounds for GP regression that complement our approach. The Supplementary Material provides a structural comparison of GP-augmented safety methods (GP--CBF Structural Comparison table).

\subsection{Policy-Agnostic Safety and Learning-Based Control}

RL has been applied to process control including voltage regulation and economic dispatch~\cite{zhang2021deep,spielberg2021dynamic}. Constrained RL (CPO, PPO-Lagrangian) enforces constraints in expectation without per-step guarantees~\cite{achiam2017constrained,ray2019reward}. Shielding~\cite{alshiekh2018safe} filters policy actions through an external safety layer. Our decoupled design is structurally equivalent to post-hoc shielding but the shield (Robust HOCBF QP) is model-based with a formal safety certificate. The policy-agnostic property transfers the guarantee to any upstream controller (LQR, MPC, PID). PPO is used as one example to demonstrate compatibility; the safety contribution comes from the Robust HOCBF filter, not from the RL algorithm.
